\BeleskeID{08-s050}
% element: 08-s050:e04
% element: 08-s050:notes
Za vreme uspona u isti izraz odziva unose se tačke u kojima je ostvareno 10\% i 90\% ukupne promene:
\begin{align*}
u_o(t_1)&=10\%\,(U-U_0)+U_0=(U-U_0)\left(1-e^{-\frac{t_1-t_0}{\tau}}\right)+U_0,\\
u_o(t_2)&=90\%\,(U-U_0)+U_0=(U-U_0)\left(1-e^{-\frac{t_2-t_0}{\tau}}\right)+U_0.
\end{align*}
Posle skraćivanja amplitude ostaje $e^{-(t_1-t_0)/\tau}=0.9$ i $e^{-(t_2-t_0)/\tau}=0.1$. Logaritmovanjem se dobijaju vremena sa slajda; njihova razlika je $\tau\ln9$.

Za kašnjenje se na isti način uzima tačka 50\%:
\[u_o(t_1)=50\%\,(U-U_0)+U_0=(U-U_0)\left(1-e^{-\frac{t_1-t_0}{\tau}}\right)+U_0.\]
Tada je eksponencijalni faktor $0.5$, pa je proteklo vreme $\tau\ln2$. Za linearno kolo prvog reda dobijaju se simetrični rezultati pri silaznoj ivici; isto važi i za odgovarajući odziv sa jednom induktivnošću.
